We present RadFI, an algebraic fault injection operator for the
empirical validation of filesystem resilience theorems. RadFI is the
falsification counterpart to the BEAMFS recovery
calculus~\cite{desbrieres2026beamfsv1}: where BEAMFS defines a
recovery operator \(\mathcal{G}\) that drives a corrupted filesystem
state back to consistency or to fail-closed, RadFI defines a
perturbation operator \(\varepsilon\) that injects controlled
bit-flips into the bio-layer payload of an instrumented filesystem,
modelling the action of a Single-Event Upset on volatile memory
transit between user space and storage.

We define \(\varepsilon\) formally over the same state space
\(\mathcal{S}\) used by BEAMFS, formulate a faithfulness property as
the statistical indistinguishability of \(\varepsilon\) from the
canonical SEU model under stated hypotheses, and derive a Linux
kernel implementation path through kprobe-based runtime
instrumentation. Three families of perturbation are treated:
write-side bit-flips on bio payloads bound for storage, read-side
bit-flips on bio payloads returning from storage (DRAM-bidirectional
model), and targeted block-level injection through a
sector-granularity filter.

This Technical Report v1 reports two empirical results obtained with
RadFI v0.1.0 to v0.1.2 against the BEAMFS implementation lineage.
First, on 2026-04-28, RadFI v0.1.0 produced a controlled bit-flip
pattern that the BEAMFS v1 recovery operator failed to correct,
exhibiting a state \(s\) such that
\(\mathcal{G}(\varepsilon(s)) \not\models\) and
\(\mathcal{G}(\varepsilon(s)) \neq \bot\). This is a counter-example
to the soundness theorem of~\cite{desbrieres2026beamfsv1} under the
actual perturbation distribution implementable by RadFI, which led
the BEAMFS author to retract Theorem IV.1 and revise the threat
model from radiation-only (v1) to full electromagnetic resilience
(v2 EMR). Second, on 2026-04-29, byte-level disk corruption of a
BEAMFS v2 INLINE-formatted filesystem
(\(\mathrm{RS}(255,239) \times 16\) per-block protection scheme)
was successfully recovered by the kernel module's
\texttt{read\_folio} path, with autonomic in-place repair persisting
the correction byte-perfect to disk. This is corroboration, not
proof, of the revised v2 recovery properties pending formal Theorem
v2.1 in~\cite{desbrieres2026beamfsv2}.

The reference implementation source is held in a private repository
(\texttt{roastercode/radfi}) pending coordinated public release
with this methodology paper. A reproducibility statement is provided in
Section~\ref{sec:repro}.
